\begin{thema}{Δ}
Στο παρακάτω σχήμα δίνεται η γραφική παράσταση της \textbf{παραγώγου} $f'$ μιας συνεχούς συνάρτησης $f$ ορισμένης στο $[-1, 7]$. Η $C_{f'}$ αποτελείται από τρία μέρη: δύο πάνω και ένα κάτω. Η $f'$ μηδενίζεται στα $x = 1$ και $x = 5$, ενώ στο $x=3$ εμφανίζει ελάχιστη τιμή. Τα επιμέρους εμβαδά είναι: $E_1 = 2$, $E_2 = 5$, $E_3 = 4$.

\begin{center}
\begin{tikzpicture}[scale=0.75]
  % grid
  \draw[help lines, gray!30] (-2,-4) grid (8,4);
  % axes
  \draw[-latex] (-2.3,0) -- (8.3,0) node[right] {$x$};
  \draw[-latex] (0,-4.3) -- (0,4.3) node[above] {$y$};
  % ticks
  \foreach \x in {-1,1,2,3,4,5,6,7}
    \draw (\x,0.1) -- (\x,-0.1) node[below,font=\footnotesize] {$\x$};
  \foreach \y in {-3,-2,-1,1,2,3}
    \draw (0.1,\y) -- (-0.1,\y) node[left,font=\footnotesize] {$\y$};
  % f' curve: positive lobe, negative lobe, positive lobe
  % Lobe 1: [-1,1], above, peak at 0
  \fill[secondcolor!20, opacity=0.5] (-1,0) .. controls (-0.5,3) and (0.5,3) .. (1,0) -- cycle;
  \draw[very thick, secondcolor] (-1,0) .. controls (-0.5,3) and (0.5,3) .. (1,0);
  % Lobe 2: [1,5], below, valley at 3
  \fill[red!15, opacity=0.5] (1,0) .. controls (1.5,-3.5) and (4.5,-3.5) .. (5,0) -- cycle;
  \draw[very thick, secondcolor] (1,0) .. controls (1.5,-3.5) and (4.5,-3.5) .. (5,0);
  % Lobe 3: [5,7], above, peak at 6
  \fill[secondcolor!20, opacity=0.5] (5,0) .. controls (5.5,3) and (6.5,3) .. (7,0) -- cycle;
  \draw[very thick, secondcolor] (5,0) .. controls (5.5,3) and (6.5,3) .. (7,0);
  % labels
  \node at (0,1.5) {$E_1 = 2$};
  \node at (3,-2) {$E_2 = 5$};
  \node at (6,1.5) {$E_3 = 4$};
  % key points
  \filldraw[secondcolor] (-1,0) circle (2pt);
  \filldraw[secondcolor] (1,0) circle (2pt);
  \filldraw[secondcolor] (5,0) circle (2pt);
  \filldraw[secondcolor] (7,0) circle (2pt);
  \node[secondcolor] at (7.5,2) {$C_{f'}$};
\end{tikzpicture}
\end{center}

\begin{erwthma}
  \item Να κατασκευάσετε πίνακα μεταβολών της $f$ στο $[-1, 7]$.
  \item Αν $f(-1) = 0$, να υπολογίσετε τις τιμές $f(1)$, $f(5)$, $f(7)$. Ποιο είναι το ολικό μέγιστο και ποιο το ολικό ελάχιστο;
  \item Να αποδείξετε ότι η εξίσωση $f(x) = 0$ έχει τουλάχιστον μία ρίζα στο διάστημα $(1, 5)$.
  \item Να βρείτε το ακριβές πλήθος των ριζών της εξίσωσης $f(x) = 0$ στο διάστημα $[-1, 7]$.
\end{erwthma}
\end{thema}
